% ======================================================================
% SECTION 5: HR 9505 - Real-Time Patient-Sponsor Direct Communication
% Act of 2026
% Revision of FDA Real-Time Clinical Trials Press Announcement April
% 2026 and FDA Decentralized Clinical Trials Guidance 2024.
%
% This file consolidates instructions previously housed in the prior
% template's Patient Priority Framework dimension 2.5 (toxicity
% warnings and AE self-management), Implementation/Metrics (Domain 3
% control-in-action subset on alert response time and false-alert
% burden), Discussion (11.2.b ePRO evidence, 11.2.f author-prior
% simulation evidence, 11.3.a FDA RTCT proof-of-concept opening),
% Limitations (12.1.e HR 4505 limitation), and Conclusions (13.1
% fifth list item) into the body of HR 9505.
% ======================================================================

[CONSOLIDATED PATIENT PRIORITY FRAMEWORK DIMENSION 2.5 (formerly
Patient Priority Framework subsection 2.5). Toxicity warnings and
adverse event self-management: ePRO, ASyMS, eRAPID. Cite
\cite{rocque2025remote_symptom_monitoring}, \cite{kearney2009asyms},
\cite{maguire2017esmart}, and \cite{absolom2021erapid}.]

[CONSOLIDATED INTRODUCTION TIE-IN AND RTCT FRAMING (formerly
Introduction block 1 closing paragraph plus Introduction block 2
item 2.i). The FDA Real-Time Clinical Trials Press Announcement of
April 28, 2026 and FDA Digital Health Technologies for Drug
Development (2026) define the regulatory opening that HR 9505 builds
on. The patient priority thesis does NOT critique the FDA; it builds
opportunistically on the FDA's own April 28, 2026 RTCT announcement
at new-trial/national-24-7-trial/FDA-April-2026/FDA_RealTime_Clinical_Trials.md
and the FDA's 2024 Decentralized Clinical Trials Final Guidance.
The FDA has validated a real-time signal-sharing conduit (Paradigm
Health) and named two proof-of-concept trials (TRAVERSE and
STREAM-SCLC). HR 9505 extends this opening into a patient-initiated
direct sponsor communication right. Cite
\cite{fda_rtct_press_2026}, \cite{fda_dct_guidance_2024}, and
\cite{fda_dhts_drug_development_2026}.]

[PROCESSING INSTRUCTION (HR 9505 - Real-Time Patient-Sponsor Direct
Communication Act of 2026, populate to 9 to 12 paragraphs structured
into the five legislative-act subsections):

The bill must explicitly state in its first paragraph that it is a
"Revision of the FDA Real-Time Clinical Trials Press Announcement of
April 28 2026 and FDA Decentralized Clinical Trials Final Guidance
of 2024" so that readers can trace the chain of prior authority.
Cite \cite{fda_rtct_press_2026}, \cite{fda_dct_guidance_2024}, and
\cite{fda_dhts_drug_development_2026}. Frame the FDA respectfully:
this bill builds on the FDA's own announced openings, not against
them.

Subsection 5.1 - Findings and Declarations: Read in full
new-trial/national-24-7-trial/FDA-April-2026/FDA_RealTime_Clinical_Trials.md
and quote FDA Commissioner Marty Makary's exact framing ("For 60
years, we've been conducting clinical trials in the same way") and
Chief AI Officer Jeremy Walsh's quote on the transformative
potential of real-time trials. Identify the two announced
proof-of-concept trials by name (AstraZeneca's TRAVERSE Phase 2 in
MCL at MD Anderson and Penn; Amgen's STREAM-SCLC Phase 1b at Mass
General). Identify Paradigm Health as the validated technical
conduit for real-time signal sharing. Compose 5 to 7 enumerated
findings stating that the FDA pilot is sponsor-initiated and that
this bill extends the same real-time channel to patient-initiated
communication.

Subsection 5.2 - Definitions:
- "Patient-sponsor direct channel" means a HIPAA-compliant, FHIR-
  based, Paradigm-Health-style aggregator endpoint that the patient
  may use to send their own structured data, modification requests,
  and AE reports directly to the trial sponsor's safety
  pharmacovigilance system. Reference
  sponsor/final_paper/scripts/sponsor_server/ (the FastAPI sponsor
  control server) and sponsor/paper/sections/safety_pharmacovigilance.tex
  for the operational baseline. Cite \cite{kawchak_2026_19396256}.
- "Real-time" means data from the patient's qualified Physical AI
  executor, ePRO endpoint, or wearable reaches the sponsor within
  one hour of generation. Reference the hourly commit cadence
  documented at new-trial/national-24-7-trial/paper/sections/results.tex
  (Simulation 1: 24 commits per day, 84 commits across 84 hours)
  and sponsor/final_paper/168_hours/ (Simulation 4: 168 GitHub
  commits across 7 days).
- "Sponsor-side patient interface" means the patient-facing endpoint
  that mirrors sponsor/final_paper/scripts/sponsor_server/ but is
  accessible to the patient under HIPAA-compliant authentication.

Subsection 5.3 - Patient-Sponsor Direct Communication Rights
(operative core):
- Right to send structured AE reports (Common Toxicity Criteria
  graded) directly to the sponsor's pharmacovigilance system within
  1 hour, drawing on the ePRO infrastructure documented at
  rocque2025remote_symptom_monitoring \cite{rocque2025remote_symptom_monitoring}
  and the ASyMS / eRAPID systems. Cite \cite{kearney2009asyms},
  \cite{maguire2017esmart}, and \cite{absolom2021erapid}.
- Right to receive sponsor-side acknowledgment of every patient
  submission within 1 hour, with the acknowledgment carrying a
  hash-chained provenance record per the MCP-PAI standard
  documented at national-platform/national_mcp/ (read all four
  chunks).
- Right to receive a per-cycle digital twin update reflecting the
  patient's own data submissions; reference
  patient-journey/stage_03_digital_twin.py (the digital-twin
  initialization Python module) and digital-twins/patient-modeling/
  for the technical baseline.
- Right to escalate any sponsor-side delay beyond 1 hour to the FDA
  RTCT pilot oversight team, mirroring the escalation pattern at
  sponsor/final_paper/scripts/coordination/ (the agent event bus and
  escalation modules).
- Right to receive a real-time read-back of how the patient's data
  is being used in the sponsor's safety analyses, with the privacy
  envelope governed by 45 CFR 46 and the data-protection framework
  in HR 9507. Cite \cite{hhs_45cfr46_protection_human_subjects}.
- Right to opt out of any individual sponsor-data exchange while
  remaining enrolled in the trial; the patient retains discretion
  over which fields the sponsor sees.

Subsection 5.4 - Implementation: 24-month deadline for ALL FDA RTCT
pilot trials to expose a patient-sponsor direct channel; FDA, ONC,
and HHS roles; explicit alignment with the FDA's own RTCT pilot
program and the FDA Oncology Center of Excellence DCT pathway. Cite
\cite{fda_oce_advancing_oncology_dct_2024} and
\cite{nih_par25_170_dht_endpoints} (the NIH funding opportunity for
DHT-derived endpoints in cancer trials).

Subsection 5.5 - Reporting and Enforcement: Annual federal report
on patient-sponsor channel uptake and median latency; civil
penalties for sponsors that fail to acknowledge within 1 hour; the
audit-trail provenance model from
sponsor/final_paper/scripts/core_agents/audit_trail_manager.py
provides the technical baseline.]

[CONSOLIDATED IMPLEMENTATION-METRICS CONTROL-IN-ACTION SUBSET
(formerly Implementation Metrics subsection 10.2 Domain 3 subset on
alert response time and false-alert burden). The federal Patient
Control Dashboard must publish, for HR 9505, the following
control-in-action metrics specifically tied to the
patient-sponsor channel: alert response time, false-alert burden,
median sponsor acknowledgment latency, and serious adverse events
detected via patient-initiated channel versus sponsor-initiated
channel. Cite \cite{rocque2025remote_symptom_monitoring}.]

[CONSOLIDATED DISCUSSION COMPARISON (formerly Discussion items
11.2.b ePRO evidence, 11.2.f author-prior simulation evidence, and
11.3.a FDA RTCT proof-of-concept opening). The ePRO and remote
symptom monitoring evidence supports HR 9505 without overclaiming.
The author's prior simulation evidence (the four LLM simulations in
new-trial/national-24-7-trial/paper/full-paper/, the single-patient
journey in patient-journey/, the 24-hour and 168-hour sponsor
simulations in sponsor/final_paper/) supplies the operational
substrate that the bill regulates. The FDA's "demonstrated
proof-of-concept" framing (TRAVERSE, STREAM-SCLC, Paradigm Health)
is the regulatory opening that HR 9505 builds on. Cite
\cite{kawchak_2026_19994945}, \cite{kawchak_2026_19119939},
\cite{kawchak_2026_19396256}, \cite{fda_rtct_press_2026}, and
\cite{fda_dhts_drug_development_2026}.]

[CONSOLIDATED LIMITATIONS PARAGRAPH (formerly Limitations item
12.1.e). HR 9505 assumes Paradigm-Health-style aggregator
infrastructure generalizes beyond TRAVERSE and STREAM-SCLC. The FDA
RTCT pilot has not yet released aggregate metrics; the 24-month
deadline is aggressive. Cite \cite{fda_rtct_press_2026}.]

[CONSOLIDATED CONCLUSIONS RESTATEMENT (formerly Conclusions item
13.1 fifth list item). HR 9505 builds respectfully on the FDA's
own April 28, 2026 RTCT framework and the 168-hour autonomous
sponsor simulation, supplying the legal infrastructure that
extends a sponsor-initiated proof-of-concept into a
patient-initiated right. Cite \cite{fda_rtct_press_2026} and
\cite{kawchak_2026_19396256}.]

Close the bill (1 paragraph after Subsection 5.5) with two
patient-control framings:

(i) FOR INDIVIDUAL PATIENTS: A 58-year-old female NSCLC patient
    matching PAT-2026-0042 sends a Grade 2 fatigue AE report from
    her home wearable on day 14 of cycle 22 of her 35 pembrolizumab
    cycles, with the sponsor's safety pharmacovigilance system
    acknowledging within 12 minutes and updating her per-cycle
    digital twin within the same hour. The patient reviews the
    read-back showing how her AE data informed the cumulative
    safety analysis at
    sponsor/final_paper/scripts/output/sponsor_24h_summary.json.

(ii) FOR BROAD ADOPTION: All FDA RTCT pilot trials must expose a
     patient-sponsor direct channel within 24 months. The bill
     positions the United States as Number 1 in the world for
     patient-initiated real-time trial communication, building
     respectfully on the FDA's own April 28, 2026 RTCT framework
     and the 168-hour autonomous sponsor simulation in
     sponsor/final_paper/168_hours/. Cite
     \cite{kawchak_2026_19396256},
     \cite{kawchak_2026_19994945}, and \cite{fda_rtct_press_2026}.
